\section{The File It Was Waiting For}
\label{sec:ch08-the-file-it-was-waiting-for}

\index{config.yaml@\code{config.yaml}!the whole of it|(}%
Here is the file, and it is the only one this chapter adds to the project:

\begin{minted}{yaml}
version: "1"

# Without this the router asks Cosmo Cloud for its execution config and
# refuses to start, because it has no token to ask with. The path is
# resolved against the directory the router is started in, not against
# this file, so start it from this folder:
#
#   router -config config.yaml
execution_config:
  file:
    path: ../graph/router.json
\end{minted}

\index{config.yaml@\code{config.yaml}!is optional}%
Strictly, even that is more than the router needs. Set
\code{EXECUTION\_CONFIG\_FILE\_PATH} in the environment and it starts and serves
the same graph with no configuration file anywhere. I write the file because a
file is a thing you edit once and keep, and an environment variable is a thing
you retype in every shell you open. The choice is yours and the router does not
care.

\index{config.yaml@\code{config.yaml}!path is relative to the working directory}%
The comment in the middle of that file earns its place. The path is resolved
against the directory the router is started in, not against the location of
\code{config.yaml}. Start the router from the folder holding the file and it
serves; start it one level up and point at \code{router/config.yaml}, byte for
byte the same file, and it does not.

\subsection{The Router Validates This File and Says So}

\index{config.yaml@\code{config.yaml}!is checked against a json schema}%
Put a key in there that does not exist and the router refuses to start:

\begin{minted}[fontsize=\footnotesize,breakanywhere]{text}
Could not load config: errors while loading config files: router config validation error for config.yaml: jsonschema validation failed with 'https://raw.githubusercontent.com/wundergraph/cosmo/main/router/pkg/config/config.schema.json#'
- at '': additional properties 'not_a_real_key' not allowed
\end{minted}

\index{config.yaml@\code{config.yaml}!the refusal names the schema by url}%
Two things about that. It is not JSON and has no level on it, unlike every
other line the router prints, because the config is validated before the
structured logger exists. And it names the schema it validated against, by url.

\index{router execution config!has no schema, unlike config.yaml}%
Put that next to what chapter~\ref{ch:composition} found one file away. The
file the composer writes has no published schema at all and is defined only by
a protobuf message in a source tree. The file you write by hand has a
JSON schema, published, and the router checks against it and hands you the
address when you fail. The generated artifact is the undocumented one. I would
have guessed the other way round.

\index{config.yaml@\code{config.yaml}!a typo cannot sit there doing nothing}%
The practical consequence is the one you want: a mistyped key in this file
cannot sit there quietly doing nothing. There is no silent-default failure mode
to debug at three in the morning, because there is no silence.

\subsection{Starting It}

\index{Cosmo Router!the startup log, read line by line}%
Compose first, then start the router from the folder holding
\code{config.yaml}:

\begin{minted}{text}
wgc router compose -i graph.yaml -o router.json
router -config config.yaml
\end{minted}

\begin{minted}[fontsize=\footnotesize,breakanywhere]{text}
INFO  Config file watching is disabled, you can still trigger reloads by sending SIGHUP to the router process
INFO  Config file provided. Values in the config file have higher priority than environment variables
      config_file=["config.yaml"]
WARN  GOMEMLIMIT was not set. Please set it manually to around 90% of the available memory to prevent OOM kills
      error="failed to set GOMEMLIMIT: cgroups is not supported on this system"
WARN  No graph token provided. The following Cosmo Cloud features are disabled. Not recommended for Production.
      features=["Schema Usage Tracking","Persistent operations","Advanced Request Tracing","Cosmo Cloud Tracing","Cosmo Cloud Metrics"]
WARN  Net poller is only available on Linux and MacOS. Falling back to less efficient connection handling method.
      error="epoll/kqueue is not supported on this system"
INFO  Prometheus metrics enabled
      listen_addr="127.0.0.1:8088" endpoint="/metrics"
INFO  Serving GraphQL playground
      url="http://localhost:3002/"
INFO  Server initialized and ready to serve requests
      listen_addr="localhost:3002" playground=true introspection=true config_version="00000000-0000-0000-0000-000000000000"
INFO  Static execution config provided. Polling and watching is disabled. Updating execution config is only possible by restarting the router
INFO  Router started
      component="supervisor"
\end{minted}

\index{Cosmo Router!six lines repeat, four are new}%
Six of those lines are ones the failed run printed too, unchanged, including
the same five disabled Cosmo Cloud features: nothing about supplying the config
locally bought any of them back. Four are new, and three of the four are worth
reading.

\index{Cosmo Router!stops polling once handed a file}%
\enquote{Polling and watching is disabled.} Handed a file, the router stops being
a client of anything. It made no network call to start and it will make none
while it runs, and the same sentence tells you the price: the config is read
once, so changing the graph means restarting the process. A router that polls a
control plane can be updated without one. This one cannot, and for a book that
is the right trade, because every state it serves is a state you produced.

\index{config\_version@\code{config\_version}!is the nil uuid chapter 7 found}%
Then look at \code{config\_version}. Thirty-two zeroes.
Chapter~\ref{ch:composition} opened \code{router.json}, found a \code{version}
field holding a nil uuid, and said it could not tell you what that field was
for. This is what it is for. The router reads it back out and logs it as the
version of the graph it is serving, which means that in a deployment fetching
configs from a control plane, this line is how you tell which one is loaded.
Composed locally, it is thirty-two zeroes on every machine, every time, and
carries nothing. Chapter~\ref{ch:blame-routing} needs the version of that
sentence where the number is real.

\index{Cosmo Router!starts with the subgraph stopped}%
One last thing about that log, and it is easiest to see by making it happen.
Stop the Sessions service and start the router: it reaches \code{Router
started} exactly as above and answers an introspection query. Nothing on 5001,
and the router does not mind. Chapter~\ref{ch:composition} made this point about
the composer, that composition reads schema files and never calls a service, and
it turns out to be true of the router at boot as well. Neither of the two things
standing between a client and your data checks that your data is there.
\index{config.yaml@\code{config.yaml}!the whole of it|)}%
